\section{Fused CPU-GPU Computing on RISC-V: A Unified RVV Architecture
for Edge and Embedded
Platforms}\label{fused-cpu-gpu-computing-on-risc-v-a-unified-rvv-architecture-for-edge-and-embedded-platforms}

\textbf{Author:} George Karani \textbf{Affiliation:} FrontForum Focus
\textbf{Email:}
\href{mailto:georgekarani44@gmail.com}{\nolinkurl{georgekarani44@gmail.com}}

\textbf{Workshop:} PARMA-DITAM 2027 \textbf{Co-located with HiPEAC,
Glasgow, Scotland, 18/01/2027}

\begin{center}\rule{0.5\linewidth}{0.5pt}\end{center}

\subsection{Abstract}\label{abstract}

Discrete GPUs provide high throughput for graphics and machine-learning
workloads, but they impose separate instruction sets, memory spaces, and
host-device data movement that are increasingly costly for
energy-constrained edge and embedded platforms. This paper presents a
fused CPU-GPU architecture built on RISC-V and its Vector (RVV)
extension, in which general-purpose RISC-V cores equipped with wide
vector units replace a discrete graphics processor. By unifying the
instruction set, register-file conventions, and memory space, the design
removes host-device transfers, simplifies the software stack, and
enables fine-grained, energy-aware scheduling across heterogeneous
workloads. We describe the microarchitecture, including a Tomasulo-style
reorder buffer, per-thread L1 caches, sparsity-aware execution, and
128-bit packed INT8/INT4 vector units, and report on a Verilog prototype
(tiny-gpu) targeting an AMD Xilinx Artix-7 FPGA carrier board.
Preliminary simulation results demonstrate correct execution of matrix
addition, matrix multiplication, and vector DP4A kernels with
configurable thread counts. We discuss trade-offs, current limitations,
and a roadmap toward cache coherence, branch divergence, and IEEE-754
FP32 support.

\textbf{Keywords:} RISC-V; RVV; unified CPU-GPU; heterogeneous
computing; vector SIMD; open-source hardware; edge AI; energy
efficiency.

\begin{center}\rule{0.5\linewidth}{0.5pt}\end{center}

\subsection{1. Introduction}\label{1-introduction}

The rapid growth of on-device artificial intelligence, real-time
graphics, and sensor processing has pushed computation toward the edge.
In these settings, the classical split between a general-purpose CPU and
a discrete GPU is increasingly inefficient: data must be copied across
buses, two distinct ISAs and programming models must be maintained, and
the energy cost of moving data often exceeds the cost of computing on
it. At the same time, RISC-V has emerged as an open, extensible, and
royalty-free ISA, and its Vector (RVV) extension provides a
standardized, variable-length SIMD programming model that is well suited
to both CPU-style control flow and GPU-style data parallelism.

This paper proposes a \textbf{fused CPU-GPU architecture} in which
RISC-V cores with RVV vector units subsume the role of a separate GPU.
Rather than treating the accelerator as a distinct device, every core is
a first-class RISC-V processor that can execute scalar control code and
wide vector kernels interchangeably. The result is a single ISA, a
unified memory space, and a single programming model, while retaining
the parallelism and throughput expected of a GPU.

Our contributions are:

\begin{enumerate}
\def\labelenumi{\arabic{enumi}.}
\tightlist
\item
  A microarchitectural blueprint for a unified CPU-GPU RISC-V many-core,
  including scheduling, memory, and execution policies.
\item
  A concrete, open-source Verilog prototype (tiny-gpu) that implements
  RV32IM, custom INT4/FP32 and 128-bit vector extensions, sparsity-aware
  execution, out-of-order execution via a Tomasulo-style reorder buffer
  (ROB), per-thread L1 caches, and a graphics pipeline.
\item
  A custom Artix-7 FPGA carrier board specification for physical
  bring-up.
\item
  Preliminary simulation results validating matrix addition, matrix
  multiplication, and vector DP4A kernels.
\end{enumerate}

The remainder of this paper is organized as follows. Section 2 reviews
background and related work. Section 3 presents the proposed
architecture. Section 4 describes the implementation. Section 5 reports
evaluation. Section 6 discusses limitations and future work. Section 7
concludes.

\subsection{2. Background and
Motivation}\label{2-background-and-motivation}

\subsubsection{2.1 The Cost of Discrete
Acceleration}\label{21-the-cost-of-discrete-acceleration}

In a conventional system, the CPU and GPU maintain separate virtual
address spaces and communicate through explicit copies over PCIe or a
shared interconnect. For edge devices, this separation manifests as
three penalties:

\begin{itemize}
\tightlist
\item
  \textbf{Data-movement energy:} transferring operands between host and
  device can dominate energy consumption for small kernels.
\item
  \textbf{Software complexity:} two ISAs, two compilers, and two
  runtimes must be maintained and synchronized.
\item
  \textbf{Latency:} kernel launch and synchronization overheads penalize
  fine-grained, interactive workloads.
\end{itemize}

\subsubsection{2.2 RISC-V and RVV as a Unifying
Substrate}\label{22-risc-v-and-rvv-as-a-unifying-substrate}

RISC-V\textquotesingle s modularity allows the same base ISA to scale
from microcontrollers to high-performance cores. The RVV extension adds
vector registers and instructions that operate on variable-length
vectors, enabling a single binary to adapt to different hardware widths.
Crucially, RVV lives in the same register file and memory model as the
scalar base ISA, so a core can interleave scalar control flow with
vector data-parallel operations without device boundaries. This makes
RISC-V + RVV a natural substrate for fusing CPU and GPU functionality.

\subsubsection{2.3 Related Work}\label{23-related-work}

Open-source GPU implementations such as Miaow and VeriGPU demonstrate
that GPU microarchitecture can be studied and extended in the open, but
they typically retain a distinct accelerator model and a complex
codebase. tiny-gpu takes a different approach: it starts from a minimal,
fully documented RISC-V scalar core and incrementally adds GPU-relevant
features (custom ML extensions, sparsity, out-of-order execution,
caches, graphics, and vector SIMD), making the design transparent and
educational while remaining extensible.

\subsection{3. Proposed Unified
Architecture}\label{3-proposed-unified-architecture}

\subsubsection{3.1 Design Principles}\label{31-design-principles}

The fused architecture follows four principles:

\begin{enumerate}
\def\labelenumi{\arabic{enumi}.}
\tightlist
\item
  \textbf{One ISA, one memory space.} All cores execute RISC-V; vector
  operations use RVV-style packed registers. There is no separate device
  address space.
\item
  \textbf{Composable cores.} Each core is a RISC-V CPU with a vector
  unit; cores are replicated to scale throughput like a GPU.
\item
  \textbf{Energy-aware execution.} Sparsity detection skips work for
  zero operands; caches reduce off-chip traffic.
\item
  \textbf{Open and sovereign.} The design is implemented in open Verilog
  and targets open toolchains, aligning with sovereign-hardware goals.
\end{enumerate}

\subsubsection{3.2 Core
Microarchitecture}\label{32-core-microarchitecture}

Each core contains:

\begin{itemize}
\tightlist
\item
  A 7-stage pipeline (IDLE, FETCH, DECODE, REQUEST, WAIT, EXECUTE,
  UPDATE).
\item
  Scalar RV32IM ALU and register file (32 x 32-bit).
\item
  Vector register file (16 x 128-bit) and vector ALU.
\item
  Per-thread private L1 data cache (2-way set-associative,
  write-through).
\item
  Tomasulo-style ROB for out-of-order writeback and hazard management.
\item
  Program counter and load-store unit.
\end{itemize}

Threads are grouped into blocks and dispatched to cores by a central
dispatcher. Each thread maintains private architectural state (PC,
scalar and vector registers, L1 cache) while executing the same kernel
in SIMD fashion.

\subsubsection{3.3 Scheduling and
Sparsity}\label{33-scheduling-and-sparsity}

A single scheduler advances a block through the pipeline. Before
EXECUTE, it checks the ROB for unresolved writes to the source
registers; if a hazard exists, the pipeline stalls until resolution or
forwarding. For ML workloads, the scheduler detects when all threads in
a block have zero-valued operands for an ALU operation and skips
EXECUTE, writing zero directly to save power and cycles.

\subsubsection{3.4 Memory System}\label{34-memory-system}

Global memory is shared and unified. Each core\textquotesingle s threads
access it through private L1 caches (write-through, LRU replacement). A
memory controller arbitrates among channels to throttle traffic and
relay responses. The data bus is 128 bits wide when vector mode is
enabled, allowing a single load/store to move a full vector register.

\subsubsection{3.5 ISA Extensions}\label{35-isa-extensions}

The prototype extends RV32IM with custom opcodes:

\begin{longtable}[]{@{}lll@{}}
\toprule\noalign{}
Opcode & Name & Description \\
\midrule\noalign{}
\endhead
\bottomrule\noalign{}
\endlastfoot
CUSTOM0 & RET & Thread retirement \\
CUSTOM1 & DP4A / FP32 & Scalar INT4 dot-product accumulate; FP32
stubs \\
CUSTOM2 & Vector SIMD & 128-bit packed vector operations \\
\end{longtable}

Vector operations include 16-lane INT8 add/multiply/MAC, 4-lane INT4
DP4A (each 8 x INT4), and FP32 stubs, plus a prefetch hint. Special
registers (\%blockIdx, \%blockDim, \%threadIdx) are mirrored in both
scalar and vector register files to support the SIMT programming model.

\subsection{4. Implementation}\label{4-implementation}

The prototype is implemented in approximately 20 SystemVerilog files and
evolved through seven milestones:

\begin{enumerate}
\def\labelenumi{\arabic{enumi}.}
\tightlist
\item
  RV32IM scalar core.
\item
  INT4/FP32 custom extensions for ML inference.
\item
  Sparsity-aware execution.
\item
  Out-of-order execution with Tomasulo-style ROB.
\item
  Per-thread L1 data cache.
\item
  Graphics pipeline (rasterizer, texture unit, framebuffer).
\item
  128-bit vector extensions (INT8/INT4 SIMD, FP32 stubs).
\end{enumerate}

\subsubsection{4.1 Hardware Platform}\label{41-hardware-platform}

For physical bring-up, a custom carrier board targets:

\begin{itemize}
\tightlist
\item
  FPGA: AMD Xilinx Artix-7 (XC7A100T-2FGG484I).
\item
  Memory: DDR3L x16 (MT41K256M16TW-107:P).
\item
  Video output: HDMI 1.4 via SII9022A transmitter.
\item
  Power: USB-C PD with multi-rail regulation (1.0 V, 1.8 V, 1.35 V, 3.3
  V).
\end{itemize}

The board was developed with a code-first workflow (Solder CLI) and is
KiCad-compatible.

\subsubsection{4.2 Verification}\label{42-verification}

Simulation uses Icarus Verilog and cocotb. Kernels are expressed as
RISC-V instruction sequences loaded into program memory; data is loaded
into unified data memory; thread count is set in the device control
register; execution is launched via a start signal.

\subsection{5. Evaluation}\label{5-evaluation}

\subsubsection{5.1 Functional
Validation}\label{51-functional-validation}

The prototype has been validated in simulation on three kernels:

\begin{itemize}
\tightlist
\item
  \textbf{Matrix addition (RV32I):} 8 threads, each computing C{[}i{]} =
  A{[}i{]} + B{[}i{]}, exercising load/store and the
  \%blockIdx/\%blockDim/\%threadIdx SIMT pattern.
\item
  \textbf{Matrix multiplication (RV32IM):} accumulated dot-products
  using mul/div and branching.
\item
  \textbf{Vector DP4A (CUSTOM2):} 128-bit vector SIMD using VDP4A.I4 for
  high-throughput INT4 dot-products.
\end{itemize}

\subsubsection{5.2 Design Parameters}\label{52-design-parameters}

\begin{longtable}[]{@{}lll@{}}
\toprule\noalign{}
Parameter & Default & Description \\
\midrule\noalign{}
\endhead
\bottomrule\noalign{}
\endlastfoot
NUM\_CORES & 2 & Compute cores \\
THREADS\_PER\_BLOCK & 4 & SIMD width per core \\
ROB\_DEPTH & 16 & ROB entries per core \\
L1\_SETS & 16 & L1 sets per thread \\
VECTOR\_ENABLE & 1 & 128-bit data bus (0 = 32-bit) \\
DATA\_MEM rows & 256 & Word-addressable data rows \\
PROGRAM\_MEM rows & 256 & Word-addressable instruction rows \\
\end{longtable}

\subsubsection{5.3 Discussion of
Metrics}\label{53-discussion-of-metrics}

Cycle counts are kernel- and configuration-dependent; the
matrix-addition test completes in approximately 1.33 x 10\^{}9
simulation nanoseconds at the configured clock (see results.xml). Rather
than reporting a single headline number, we emphasize correctness,
configurability, and the ability to study microarchitectural trade-offs
(cache size, ROB depth, vector width) in an open, modifiable design.
Area and power estimates will be reported after synthesis on the target
FPGA.

\subsection{6. Limitations and Future
Work}\label{6-limitations-and-future-work}

\begin{itemize}
\tightlist
\item
  \textbf{FP32 is stubbed.} Vector and scalar FP32 paths are
  placeholders pending IEEE-754-compliant IP.
\item
  \textbf{No cache coherence.} Caches are per-thread and write-through;
  a shared L2 and coherence protocol are future work.
\item
  \textbf{No branch divergence handling.} All threads in a block execute
  in lockstep.
\item
  \textbf{No multi-warp scheduling.} One block runs to completion per
  core.
\item
  \textbf{Memory coalescing} is not implemented.
\item
  \textbf{Physical bring-up} of the carrier board is in progress.
\end{itemize}

\subsection{7. Conclusion}\label{7-conclusion}

This paper presented a fused CPU-GPU architecture on RISC-V that
replaces a discrete accelerator with vector-capable RISC-V cores,
unifying ISA, memory, and programming model. An open Verilog prototype
demonstrates the approach across seven milestones, from a scalar RV32IM
core to 128-bit vector SIMD, with sparsity-aware and out-of-order
execution. Preliminary simulation validates core kernels. The design is
intended as a transparent, extensible platform for research in
heterogeneous, energy-efficient, and sovereign computing, and aligns
with the PARMA-DITAM focus on RISC-V, accelerators, and embedded
multicore platforms.

\begin{center}\rule{0.5\linewidth}{0.5pt}\end{center}

\subsection{References}\label{references}

\begin{enumerate}
\def\labelenumi{\arabic{enumi}.}
\tightlist
\item
  A. Waterman and K. Asanovic, "The RISC-V Instruction Set Manual,
  Volume I: User-Level ISA," RISC-V Foundation, 2019.
\item
  RISC-V International, "RISC-V Vector Extension (RVV) Specification,"
  2021.
\item
  K. Asanovic et al., "The Rocket Chip Generator," UC Berkeley, 2016.
\item
  M. Owaida et al., "Miaow: An Open Source GPU," 2015.
\item
  H. Perkins, "VeriGPU: An Open Source Verilog GPU," 2020.
\item
  J. L. Hennessy and D. A. Patterson, "Computer Architecture: A
  Quantitative Approach," 6th ed., Morgan Kaufmann, 2019.
\item
  G. Karani, "tiny-gpu: A RISC-V-based GPU in Verilog," GitHub, 2026.
\item
  Dagstuhl Publishing, "OASIcs - OpenAccess Series in Informatics,"
  \url{https://www.dagstuhl.de/oasics}.
\end{enumerate}

\begin{center}\rule{0.5\linewidth}{0.5pt}\end{center}

\emph{Draft prepared for PARMA-DITAM 2027. Author metadata and
references to be finalized before submission.}
